\begin{table}[t]
\centering
\scriptsize
\setlength{\tabcolsep}{4pt}

\begin{tabular}{lc|lclc}
\toprule
\multicolumn{2}{c|}{\textbf{FLEURS}} &
\multicolumn{4}{c}{\textbf{DoReCo}}\\
\midrule
\textbf{Family} & \textbf{\#} &
\textbf{Family} & \textbf{\#} &
\textbf{Family} & \textbf{\#} \\
\midrule
Indo-European & 41 &
Austronesian & 7 &
Mixed Language & 1 \\
Niger-Congo & 9 &
Sino-Tibetan & 4 &
Mixe-Zoque & 1 \\
Afroasiatic & 7 &
Indo-European & 3 &
Nilo-Saharan & 1 \\
Austronesian & 5 &
Niger-Congo & 3 &
Pama-Nyungan & 1 \\
Turkic & 5 &
Afroasiatic & 2 &
Pano-Tacanan & 1 \\
Dravidian & 4 &
Arawakan & 2 &
Trans-New Guinea & 1 \\
Sino-Tibetan & 3 &
Austroasiatic & 2 &
Tungusic & 1 \\
Uralic & 3 &
Nakh-Daghestanian & 2 &
Tuu & 1 \\
Austroasiatic & 2 &
Turkic & 2 &
Uralic & 1 \\
Kra-Dai & 2 &
Algic & 1 &
Yam & 1 \\
Japonic & 1 &
Boran & 1 & \\
Kartvelian & 1 &
Chibchan & 1 & \\
&
Gunwinyguan & 1 & \\
&
Isolate & 1 & \\
&
Kartvelian & 1 & \\
&
Koreanic & 1 & \\
&
Mayan & 1 & \\
\bottomrule
\end{tabular}

\caption{Language family distributions in the multilingual evaluation datasets. FLEURS (85 languages) is dominated by Indo-European languages, whereas DoReCo (45 languages) covers a substantially broader range of language families.}
\label{tab:family_stats}
\end{table}